% =====================================================================
% 13. Reflection and Critical Evaluation
% =====================================================================
\section{Reflection and Critical Evaluation}
\label{sec:reflection}

\todo{This section must be in your own voice. The points below are prompts drawn from how the project has developed; keep the ones that are true, add your own, and rewrite them in the first person.}

\begin{itemize}
    \item \textbf{Industry work versus research.} The first plan followed the partner's rollout schedule and had no time for the research experiments.
    Separating the two sets of deliverables was the most important planning decision so far.
    \item \textbf{Dropping the user study.} Removing human participants simplified ethics and scheduling but weakened the evidence available for RQ4. The response was to combine three analytical measures and to state the limits plainly rather than overstate them.
    \item \textbf{Designing for blocked access.} Assuming rig access would arrive on time was risky.
    The adapter boundary and simulator turned a likely blocker into a known fallback.
    \item \textbf{Assumptions to revisit:} the provisional performance targets, the choice of time-series store, and whether ISA-18.2 benchmarks are appropriate for an intermittently watched test rig.
\end{itemize}
